%% kapa: exam=1 round=33 year=2022 period=1 form=A subject=civil_law subject_ko=민법 seq=3 number=3
%% source: Q-Net 감정평가사 기출문제 (cst00309) artlSeq=5209997 file=raw/1008/5209997/2022년도 제33회 감정평가사 1차시험 1교시(민법,경제학,부동산학).pdf
%% conversion: kordoc 4.24.51 -> tools/gen1.py (auto); review=proofread-sonnet; flags=none
\begin{question}{3}
미성년자 甲과 그의 유일한 법정대리인인 乙에 관한 설명으로 옳은 것은? (다툼이 있으면 판례에 따름)
\begin{choices}
  \item 甲이 그 소유 물건에 대한 매매계약을 체결한 후에 미성년인 상태에서 매매대금의 이행을 청구하여 대금을 모두 지급받았다면 乙은 그 매매계약을 취소할 수 없다.
  \item 乙이 甲에게 특정한 영업에 관한 허락을 한 경우에도 乙은 그 영업에 관하여 여전히 甲을 대리할 수 있다.
  \item 甲이 乙의 동의 없이 타인의 적법한 대리인으로서 법률행위를 했더라도 乙은 甲의 제한능력을 이유로 그 법률행위를 취소할 수 있다.
  \item 甲이 乙의 동의 없이 신용구매계약을 체결한 이후에 乙의 동의 없음을 이유로 그 계약을 취소하는 것은 신의칙에 반한다.
  \item 乙이 재산의 범위를 정하여 甲에게 처분을 허락한 경우, 甲이 그에 관한 법률행위를 하기 전에는 乙은 그 허락을 취소할 수 있다.
\end{choices}
\end{question}
\answer{5}
